\section{处理器：数据通路}

\subsection{单周期 CPU 概述}

单周期 CPU 在\textbf{一个时钟周期内}完成从取指到写回的全部操作。CPI = 1，但时钟周期受最慢指令（如 $\texttt{lw}$）限制。

\subsection{数据通路组件}

\begin{table}[H]
\centering
\caption{基本数据通路组件}
\label{tab:datapath_components}
\begin{tabulary}{\textwidth}{CCC}
\toprule
\textbf{组件} & \textbf{功能} & \textbf{控制信号} \\
\midrule
程序计数器（PC） & 存放下一条指令地址 & PCWrite \\
指令存储器 & 按地址取指令 & — \\
寄存器堆 & 32 个 32 位通用寄存器 & RegWrite \\
ALU & 算术/逻辑运算 & ALUOp, ALUSrc \\
数据存储器 & 读写数据 & MemRead, MemWrite \\
多路选择器（MUX） & 选择数据来源 & 选择信号 \\
符号扩展单元 & 16-bit $\to$ 32-bit 扩展 & — \\
\bottomrule
\end{tabulary}
\end{table}

\subsection{基本指令的执行}

\subsubsection{R 型指令}

\begin{equation}
\text{add } \$rd, \$rs, \$rt \quad \longrightarrow \quad \$rd = \$rs + \$rt
\end{equation}

数据流：$\text{PC} \to \text{取指} \to \text{译码（读} \$rs, \$rt\text{）} \to \text{ALU} \to \text{写回} \$rd \to \text{PC+4}$

\subsubsection{I 型：lw}

\begin{equation}
\text{lw } \$rt, \text{offset}(\$rs) \quad \longrightarrow \quad \$rt = M[\$rs + \text{SignExt}(16\text{-bit offset})]
\end{equation}

数据流：$\text{PC} \to \text{取指} \to \text{译码（读} \$rs\text{）} \to \text{ALU 算地址} \to \text{访存} \to \text{写回} \$rt \to \text{PC+4}$

存储器访问额外需要 ALU 用于地址计算，结果经数据存储器读出。

\subsubsection{I 型：sw}

\begin{equation}
\text{sw } \$rt, \text{offset}(\$rs) \quad \longrightarrow \quad M[\$rs + \text{SignExt}(16\text{-bit offset})] = \$rt
\end{equation}

与 lw 类似，但写入方向相反——寄存器的值写入存储器而非读回。

\subsubsection{I 型：beq}

\begin{equation}
\text{beq } \$rs, \$rt, \text{label} \quad \longrightarrow \quad \text{if } \$rs = \$rt \text{ then } \text{PC} \leftarrow \text{PC}+4 + \text{offset}\times4
\end{equation}

ALU 做减法判断是否相等。若相等，PC 取分支目标（PC+4 + offset$\times$4），否则取 PC+4。

\subsection{完整单周期数据通路}

将上述组件用多路选择器和总线连接，形成完整的 MIPS 单周期数据通路。控制信号由控制器根据指令操作码生成。

\begin{itemize}
    \item \textbf{取指阶段}：PC 送指令存储器，取出 32 位指令
    \item \textbf{译码阶段}：寄存器堆读出 rs、rt 值，符号扩展立即数
    \item \textbf{执行阶段}：ALU 执行运算（算术 / 地址计算 / 比较）
    \item \textbf{访存阶段}：读/写数据存储器（仅 lw / sw）
    \item \textbf{写回阶段}：结果写回寄存器堆（ALU 结果 / 存储器数据）
\end{itemize}

\subsection{关键路径与性能}

单周期 CPU 的时钟周期由最长路径（$\texttt{lw}$ 指令：取指+译码+ALU+访存+写回）决定。这导致简单指令（如 $\texttt{add}$）也必须等待相同的长周期，效率很低。这是后续流水线设计的直接动机。
